% I. Теоретична основа, проектиране и програмна реализация.
% Теорията е около една страница: проектът се предава преди курсът да е започнал.
% Всичко, което вече е казано с число в Раздел III, тук се посочва с препратка, не се повтаря.
\section{Теоретична основа, проектиране и реализация}
\label{sec:theory}

\subsection{Метаевристики и паралелни схеми}
\label{subsec:theory-background}

Задачата за комбинаторна оптимизация се задава с крайно множество допустими решения и целева функция
върху него, но множеството расте толкова бързо с размера на входа, че изчерпателното му обхождане е
неприложимо. Метаевристиките управляват подчинена евристика, която сама би заседнала в локален
оптимум, и се различават по това как балансират изследването и
задълбочаването~\cite{blum2003metaheuristics,talbi2009metaheuristics}. В проекта влизат четири от тях,
покриващи четири различни отношения към споделеното състояние.
\term{Симулираното закаляване}~\cite{kirkpatrick1983annealing} приема влошаващи ходове с вероятност,
намаляваща с температурата, тоест зависи изцяло от схемата на охлаждане и от съседството.
\term{Търсенето с забрани}~\cite{glover1986future} пази краткосрочна памет за скорошните ходове, а
забраненият списък е структура с честа проверка за принадлежност, тоест точка, в която разположението
на данните влияе на времето за итерация. \term{Генетичният алгоритъм}~\cite{holland1975adaptation}
оценява индивидите на популацията независимо, тоест се разпределя естествено между нишки.
\term{Мравчената колония}~\cite{dorigo1996antsystem} строи решения над обща матрица от следи, тоест
споделеното състояние тук е част от алгоритъма.

Паралелните метаевристики се класифицират по това какво точно се
разпределя~\cite{crainic2010parallel,alba2005parallel}. Независимият многократен старт не изисква
комуникация и мащабира почти линейно, но не използва натрупаното от другите нишки; кооперативните
схеми го използват, но плащат със синхронизация и с риск от преждевременна сходимост; островният модел
стои между двете и въвежда честотата на обмена като допълнителен параметър. Сравнимостта с публикувани
резултати изисква инстанции с известни стойности, каквито са TSPLIB~\cite{reinelt1991tsplib} и
OR-Library~\cite{beasley1990orlibrary}, а тълкуването на измерена производителност изисква модел,
отделящ ограничението от изчисленията от ограничението от паметта; моделът на покривната
линия~\cite{williams2009roofline} дава такъв, но изисква хардуерни броячи, каквито тук не са снети.

Източниците описват стратегиите и схемите поотделно и всеки фиксира по едно сравнение. Липсва общата
рамка, в която задача, стратегия и схема за изпълнение се менят независимо при едни и същи условия, а
без нея всяко сравнение между две публикувани реализации е и сравнение между два набора инстанции и
два критерия за спиране.

\subsection{Проектни решения}
\label{subsec:theory-choices}

Четирите независими решения и основанията им са в Табл.~\ref{tab:alt-summary}
\tabref{tab:alt-summary}. Динамичният полиморфизъм е избран срещу първоначалното намерение за шаблони
по горещия път, защото оценяването заема под 11 на сто от една итерация, §\ref{subsec:res-profile}.

\begin{table}[htbp]
\centering
\renewcommand{\arraystretch}{0.9}
\caption{Проектни решения и основанието за всяко от тях}
\label{tab:alt-summary}
\begin{adjustbox}{max width=\textwidth}
\begin{tabular}{@{}L{2.4cm}L{3.4cm}L{2.0cm}L{6.2cm}@{}}
\toprule
\textbf{Решение} & \textbf{Алтернативи} & \textbf{Избрано} & \textbf{Основание} \\
\midrule
Език и среда & управлявана среда; език без стандартен паралелизъм &
C\texttt{++}20~\cite{iso_cpp20} & Контрол върху разположението на данните, без събирач на боклук в горещия цикъл. \\
Паралелизъм & задачно ориентирана среда; директиви над цикли &
явни нишки & Предметът на измерването е цената на синхронизацията и тя не бива да е скрита зад чужд планировчик. Задачната среда остава неизмерена. \\
Договор за задача & статичен полиморфизъм с шаблони &
динамичен & Оценяването е под 11 на сто от итерацията, §\ref{subsec:res-profile}. \\
Обмен между нишките & само една от трите схеми &
и трите & Образуват редица с нарастваща комуникация; кое се изплаща е проверено в §\ref{subsec:res-cooperation}, а не допуснато. \\
\bottomrule
\end{tabular}
\end{adjustbox}
\end{table}

\subsection{Архитектура, договори и реализация}
\label{subsec:design-layers}

Изискването е едно: задачата, метаевристиката и схемата за паралелно изпълнение трябва да се менят
независимо, иначе сравнението между схемите става и сравнение между две реализации на алгоритъма.
Оттук следват четири слоя със зависимости само надолу: договор за задача (\code{evaluate},
\code{delta}, \code{apply}, \code{random\_move}, \code{construct}), метаевристики (\code{step},
\code{best}, \code{inject}), схеми за паралелно изпълнение и измервателна рамка. Слоят на алгоритмите
не знае колко нишки има, а слоят на задачите не знае кой алгоритъм го вика, тоест добавянето на трета
задача не наложи промяна в нито един алгоритъм. И трите задачи използват едно представяне, вектор от
цели числа, а кандидатът носи кеширано състояние, тоест общото тегло при раницата и обратната
пермутация при търговския пътник.

Метаевристиката е обект със състояние, изпълняващ стъпки, защото нишка трябва да може да продължи от
решение, дошло отвън. Една стъпка обаче означава различно нещо при различните алгоритми, тоест
бюджетът се задава във време или в брой оценявания, никога в стъпки, а периодът на обмен се брои в
оценявания: период, разумен за отгряването с милиони итерации в секунда, означава нула обмени за
мравчената колония с около сто и петдесет, и първата версия, броила итерации, наистина не обменяше
нищо. Трите схеми се различават само по обмена: при многократния старт обмен няма и той е еталонът за
цената на комуникацията; при кооперативната нишките обновяват обща най-добра стойност; при островния
подпопулациите обменят по пръстен, а не по пълен граф, където схемата би се изродила в кооперативната.

Реализацията е около 4300 реда C\texttt{++}20 без нито една задължителна външна зависимост
\tabref{tab:files} и е публична на \url{https://github.com/Alex-Tsvetanov/anneal}. Липсата на
зависимости е условие за възпроизводимост: изграждане, което изисква мениджър на пакети, трето лице
няма да повтори. Затова изпълнителят на тестове е сто реда вместо външна рамка, а генераторът е
xoshiro256++, защото стандартните разпределения не са специфицирани до бит. Три решения са взети по
измерване, а не по очакване, и числата им са в Раздел~\ref{sec:results}: генераторът на съседен ход,
чиято първа версия правеше търсенето безполезно; изборът между матрица на разстоянията и координати;
и подравняването на данните на всяка нишка на 64 байта. Тестовете са 28 и покриват четири твърдения:
че инкременталното оценяване дава същия резултат като пълното и при трите задачи, че всяко върнато
решение е допустимо, че при фиксиран бюджет в оценявания резултатът е побитово възпроизводим, и че
инстанциите с обявен оптимум наистина го имат, защото без такъв свидетел оптимумът е само твърдение.

\begin{table}[htbp]
\centering
\renewcommand{\arraystretch}{0.9}
\caption{Модули на реализацията, брой редове от \code{wc -l}}
\label{tab:files}
\begin{adjustbox}{max width=\textwidth}
\begin{tabular}{@{}L{3.4cm}L{9.0cm}r@{}}
\toprule
\textbf{Файл} & \textbf{Съдържание} & \textbf{Редове} \\
\midrule
\code{include/anneal/} & договорите \code{Problem} и \code{Metaheuristic}, xoshiro256++, статистики, TSPLIB & 880 \\
\code{src/algorithms.cpp} & четирите метаевристики & 445 \\
\code{src/parallel.cpp} & трите схеми и споделеното състояние & 320 \\
\code{src/tsp.cpp}, \code{knapsack}, \code{coloring} & трите задачи и генераторите на инстанции & 883 \\
\code{benchmarks/}, \code{apps/}, \code{src/experiment.cpp} & измервателна програма с осем подкоманди, демонстрация, рамка & 1010 \\
\code{tests/} & изпълнител на тестове и 28 теста & 807 \\
\bottomrule
\end{tabular}
\end{adjustbox}
\end{table}
